% Generated draft: requires mathematical and source review.
\documentclass[11pt]{article}
\usepackage[margin=25mm]{geometry}
\usepackage{fontspec}
\setmainfont{Latin Modern Roman}
\usepackage{amsmath,amssymb,mathtools,unicode-math}
\setmathfont{Latin Modern Math}
\usepackage{enumitem,needspace,xurl,hyperref}
\setlength{\parindent}{0pt}
\setlength{\parskip}{5pt}
\setlength{\emergencystretch}{3em}
\newcommand{\N}{\mathbb{N}}
\newcommand{\Z}{\mathbb{Z}}
\newcommand{\Q}{\mathbb{Q}}
\newcommand{\R}{\mathbb{R}}
\newcommand{\C}{\mathbb{C}}
\newcommand{\F}{\mathbb{F}}
\newcommand{\A}{\mathbb{A}}
\newcommand{\PP}{\mathbb{P}}
\newcommand{\E}{\mathbb{E}}
\newcommand{\eps}{\varepsilon}
\newcommand{\dd}{\,\mathrm{d}}
\begin{document}
\section*{UnsolvedMath ID: 20002863}
\textbf{Source problem code:} catalogue entry
\subsection{Definitions and mathematical statement}
Let $z_0\in M\subset\C^n$, where $M$ is a real-analytic generic submanifold, meaning $T_xM+iT_xM=T_x\C^n$ at every nearby $x\in M$. Let $f$ be a meromorphic map germ to $\C^N$ at $z_0$, represented on a neighborhood $U$ by
\[
f=(p_1/q,\ldots,p_N/q),\qquad p_j,q\in\mathcal O(U),\quad q\not\equiv0.
\]
The extracted ``$f=pq$'' is interpreted as $f=p/q$, as required by the meromorphic-map context; the original typography is not reproduced in the record. Let $W\subset U\setminus M$ be a one-sided wedge attached to $M$ near $z_0$: locally it approaches an open piece of the edge $M$ through an open cone of transverse directions. Assume $f$ is holomorphic on $W$. Meylan's question is whether the boundary condition that $\|f(w)\|^2$ ``increases to $1$ as $w\to M$ in $W$'' forces a holomorphic extension to a full neighborhood of $z_0$. Holomorphic extension is the substantive target, since a meromorphic germ on a neighborhood is already assumed. The source reports that the real-codimension-one case is known, attributing it to Chiappari (1991).

The supplied partial result uses the explicitly stronger edge-limit reading
\[
\forall x\in M\cap U\quad \lim_{\substack{w\to x\\w\in W}}\|f(w)\|^2=1,
\]
where every wedge approach is allowed after shrinking $U$. This does not by itself formalize monotonicity implicit in the word ``increases''; a requirement $\|f\|<1$ on the wedge would be an additional interpretation. Under the displayed full-limit assumption the aid asserts the real-analytic boundary identity
\[
|q|^2=\sum_{j=1}^N|p_j|^2\quad\hbox{on }M\cap U.
\]
In particular, $q(x)=0$ on the edge implies $p_j(x)=0$ for every $j$.

For an irreducible local complex hypersurface $H$ in the zero divisor of $q$, call $E\subset H_{\mathrm{reg}}\cap M$ a generic uniqueness piece if it is a nonempty relatively open piece of a real-analytic submanifold of $H_{\mathrm{reg}}$ with
\[
T_xE+iT_xE=T_xH_{\mathrm{reg}},
\]
and hence vanishing of a holomorphic function on $E$ forces its restriction to the relevant irreducible hypersurface germ to vanish identically. The reported cancellation criterion is that any polar component with such an edge uniqueness piece divides every numerator and therefore cancels as a common factor. A reduced representation means no irreducible factor of $q$ divides all the $p_j$. In such a representation every surviving polar hypersurface must avoid $W$ and have no such generic edge uniqueness piece; its intersection with the edge is correspondingly degenerate. If every putative polar component has such a uniqueness piece in the required local germ sense, no reduced denominator survives and $f$ extends holomorphically near $z_0$.

This is a conditional removability criterion. It does not show that every polar component in Meylan's setting has a generic edge slice. The aid also reports a rational inner function on the bidisk whose norm tends to one radially at a singular edge point while the meromorphic singularity is nonremovable. Its actual quotient is not supplied here. That reported example concerns a weaker single radial approach and does not refute the full wedge-limit problem.
\subsection{Short English statement}
Meylan asks whether a meromorphic map, holomorphic on an attached wedge and approaching the unit sphere at the edge, must extend holomorphically across the edge. Under a full limit at every nearby edge point, the supplied aid gives a boundary identity and a cancellation criterion for polar hypersurfaces with generic edge intersections. Degenerate intersections remain outside that criterion; radial convergence alone is reported to be too weak.
\subsection{Sources}
\begin{itemize}
\item Local registry: permanent catalogue entry, exact source code catalogue entry; \texttt{aim-several-complex-variables-notes.json}, \texttt{notes[20]}.
\item AIM workshop \emph{Emerging applications of complexity for CR mappings}, Palo Alto, 9--13 August 2010, Problem 2 attributed to F. Meylan, \url{https://aimath.org/WWN/crmappings/crmappings.pdf}. The verification field identifies a ``Mappings between balls'' section but ends before reproducing the problem.
\item Original remark: the codimension-one case is known, Chiappari, 1991. No title, journal, theorem number, or complete authoritative reference is supplied; the original literature field is empty.
\item Supplied research attempt 1, registry updated 14 August 2026: partial result, low novelty confidence, machine-generated and not peer reviewed. It supplies the polar-edge uniqueness criterion and reports a bidisk radial-limit obstruction, without its explicit quotient or a full external citation in this record.
\end{itemize}
\end{document}
